% !TeX spellcheck = it_IT

\documentclass[format=169, handout]{beamer}

\usepackage[utf8]{inputenc}
\usepackage[italian]{babel}
\usepackage{multicol}

\usetheme{metropolis}

\usepackage{mathtools}
\DeclarePairedDelimiter{\ceil}{\lceil}{\rceil}

\title{Esercitazioni di Informatica B}
\subtitle{Struct}

\author{Stefano Cereda\\
	stefano.cereda@polimi.it
}
\date{30/10/2018}
\institute[PoliMi]{\vspace{0.5cm}\centering Politecnico di Milano \\ \vspace{0.2cm}
	\includegraphics[width=0.2\textwidth]{../logopolimi}}

\setbeamercovered{invisible}

\makeindex

\begin{document}
	\begin{frame}
	\maketitle
\end{frame}

\begin{frame}<handout:0>{Nuovi esercizi su argomenti vecchi}
Aggiunti due esercizi sulle matrici e alcuni esercizi sul binario.
\end{frame}

\begin{frame}[allowframebreaks]{Scontrini}
Un supermercato ha memorizzato il proprio archivio di scontrini nell'array \emph{archivio}, in cui ogni elemento è una struttura dati di tipo \emph{scontrino} con le  seguenti informazioni:
\begin{itemize}
	\item IDcliente: stringa di 10 caratteri che identifica univocamente il cliente
	\item totale: totale della spesa in Euro
	\item punti: punti premio associati alla spesa
\end{itemize}
Ad ogni spesa, vengono assegnati ad ogni cliente un quantitativo di punti premio pari alla somma dei punti raccolti più un ulteriore punto premio per ogni 10 euro spesi.

\begin{enumerate}
	\item  Definire il tipo $scontrino$ e l'array $archivio$
	\item Scrivere il frammento di codice C che, dato il codice identificativo di un cliente, calcoli la somma dei suoi punti premio.
	\item Scrivere il frammento di codice C che ordina l'array \emph{archivio} secondo il totale della spesa in ordine decrescente.
\end{enumerate}
Nota: Si faccia attenzione al fatto che un cliente può comparire in più di uno scontrino.
\end{frame}

\begin{frame}{Anagrafe}
Si vuole gestire l’anagrafe della motorizzazione e dei proprietari di auto. Un veicolo è identificato da:
\begin{itemize}
\item Targa
\item Modello
\item Proprietario (si suppone che ogni auto abbia un unico proprietario)
\end{itemize}

I proprietari di auto sono identificati dai campi:
\begin{itemize}
\item Cognome
\item Nome
\item Un insieme di indirizzi di recapito (max 3).
\end{itemize}

Un recapito è composto da: via, numero civico e CAP. L’anagrafe contiene al max N veicoli.
\end{frame}

{
\small
\begin{frame}{Anagrafe}
Definire in C la struttura dati che permette di rappresentare la situazione descritta.

Scrivere il programma C che chiede all’utente di inserire i dati degli N veicoli da memorizzare
nell’anagrafe. Per ogni veicolo il programma deve acquisire anche le informazioni relative al
proprietario, inclusi i suoi recapiti. In una prima versione del programma si inseriscano esattamente
tre recapiti per proprietario, in una seconda versione si consenta all’utente di inserire un numero di
recapiti minore o uguale a tre.

Finita la fase di inserimento il programma stampa il contenuto della struttura nel formato:\\
VEICOLO Targa XXXXX \\
PROPRIETARIO Cognome XXXXXXXXX \\
RECAPITI CAP XXXXXXXXX ... XXXXXXXXX \\
\end{frame}
}



\begin{frame}{Fantacalcio da TdE 17/09/2012}
Definire un tipo di dato \emph{fantacalcio} che contiene una lista di $N = 20$ \emph{fantasquadre}. Ogni fantasquadra è identificata dal nome, il colore della casacca, il punteggio corrente e l’allenatore. L’\emph{allenatore} è definito dal nome, cognome e numero di coppe vinte (un intero $\geq 0$). Infine si dichiari una variabile \emph{campionato} di
tipo fantacalcio. 

Scrivere un frammento di codice in linguaggio C che stampi a video il numero totale degli allenatori (contenuti all'interno della variabile \emph{campionato}) che hanno vinto almeno una coppa.
\end{frame}



\begin{frame}[allowframebreaks]{Archivio reperti da TdE 26 Giugno 2017 - vedere pdf esame}
Un’organizzazione che si occupa di archeologia sottomarina vuole sviluppare un software in linguaggio C
che permetta di catalogare tutti i reperti archeologici che sono stati scoperti. Ogni reperto è caratterizzato
dalle seguenti informazioni:
\begin{itemize}
\item coordinate geografiche in cui si trova
\item profondità del reperto rispetto al livello del mare
\item  peso del reperto
\item  tipo di reperto, questo può essere: manufatto in terracotta, manufatto in marmo, manufatto in metallo
\end{itemize}
Tutti i reperti vengono memorizzati in un array di dimensione massima 1000.

\begin{enumerate}
\item Si definiscano in linguaggio C le strutture dati e le variabili necessarie per rappresentare l’insieme dei
reperti nella variabile \emph{insiemeReperti}.
\item  Si supponga che la variabile \emph{insiemeReperti} sia stata precedentemente riempita con le informazioni
di 1000 reperti e si sviluppi in C il ciclo che, per ogni reperto, stampa a video il suo peso. Si dichiarino
anche tutte le variabili necessarie allo scopo che non sono state dichiarate nel punto precedente.

Periodicamente, l’organizzazione progetta spedizioni per raccogliere i reperti scoperti. L’operazione non è
sempre possibile perché il veicolo che viene usato per raccogliere i reperti ha una portata massima (un peso
massimo che può portare) che non può essere superato. D’altra parte, l’organizzazione vuole fare in modo
che il veicolo venga utilizzato nel modo più efficiente possibile e quindi la spedizione verrà pianificata in
modo tale che il veicolo raccolga il maggior numero possibile di reperti nell'area prescelta. La spedizione
viene pianificata come segue:

Si parte con un veicolo scarico, di portata \emph{pMax}, da una coppia di coordinate geografiche (\emph{latitudineI}
e \emph{longitudineI}) e si esplorano le vicinanze di quelle coordinate geografiche in modo da rimanere
nell'ambito della zona di ricerca, i cui confini sono definiti dai punti $latitudineI ± soglia$,
$longitudineI ± soglia$. Per ogni reperto che si trova nei confini definiti, si pianifica di caricare il
reperto a bordo del veicolo, a patto che il suo peso, sommato al peso del carico corrente portato dal veicolo,
sia minore di $pMax$. In caso contrario, quel reperto viene scartato e si procede a cercarne un altro, sempre
rimanendo nei confini dell'area di ricerca. La ricerca termina quando il peso dei reperti da trasportare avrà
raggiunto il valore $pMax$ oppure quando tutti i reperti nella zona di ricerca saranno stati esaminati.

L’algoritmo di pianificazione del viaggio, durante l’analisi sopra descritta, crea un array di reperti chiamato
$caricoVeicolo$ che contiene, alla fine dell'analisi, tutti i reperti che dovranno essere caricati sul veicolo
durante l’esecuzione della spedizione.

\item Dato un reperto che assumiamo si trovi nella cella $i$-esima (con $i$ che ha un valore qualsiasi compreso tra
0 e 999, estremi inclusi) dell'array $insiemeReperti$, si scriva la condizione che, se vera, ci permette di
dire che il reperto si trova nell'area di ricerca definita dai confini indicati sopra.

\item Si sviluppi il frammento di codice in C necessario per acquisire da tastiera: la portata massima del
veicolo, le coordinate geografiche di partenza della spedizione, il valore soglia che definisce i confini di
ricerca. Si dichiarino tutte le variabili necessarie allo scopo.

\item Si sviluppi in linguaggio C l’algoritmo di pianificazione del carico descritto in precedenza. Oltre alle
dovute inizializzazioni di variabili (non è necessario inizializzare $insiemeReperti$ che si assume essere
già riempito), la parte principale di questo algoritmo sarà costituita a un ciclo che analizza tutti i reperti in
$insiemeReperti$ e, per ogni reperto:
\begin{itemize}
\item  Valuta se si trova nell'area di ricerca (si faccia riferimento alla soluzione fornita per il punto 3)
\item  In caso positivo, valuta se può essere caricato sul veicolo
\item  In caso positivo, aggiorna la variabile $caricoVeicolo$ in modo che contenga i dati del reperto in
questione
\end{itemize}
Si dichiarino anche tutte le variabili necessarie allo scopo che non sono state dichiarate nei punti precedenti.
\end{enumerate}

\end{frame}

\end{document}